\begin{textsample}{POS Dim 1 – pb – Score 57.00 – pb/pb\_em\_74.txt}  \label{ex:f1_pos_146}
2.3.10 PROCESSOS \textbf{METODOLÓGICOS} POR ÁREA, COM AS ESPECIFICIDADES DOS COMPONENTES Com a ampliação de redes de interações sociais mediadas pela \textbf{tecnologia}, os campos de aprendizagens na atualidade não se \textbf{restringem} apenas ao espaço escolar.
O crescente avanço da computação e das \textbf{tecnologias} \textbf{digitais} provocaram modificações importantes na forma de pensar, agir, ensinar, aprender, se comunicar, e de trabalhar.
A computação e as novas \textbf{tecnologias} revolucionaram a forma que \textbf{vivemos} e, em especial, a forma de interação das pessoas.
No ambiente escolar, não foi diferente.
O espaço dedicado à aprendizagem sofreu profundas variações, e a \textbf{mobilização} de conhecimentos ocorre hoje nas mais diversificadas situações, seja em ambientes físicos ou em espaços e tempos \textbf{virtuais}.
Surge, assim, uma preocupação em aliar as múltiplas formas de produção de conhecimento, ao espaço/tempo escolar, na tentativa de resgatar o interesse dos estudantes pelo universo científico, estimulando -os a conhecer a história por trás do desenvolvimento científico e oportunizando -lhes a \textbf{mobilização} do conhecimento de forma que possam compreender, interagir e modificar cenários nos quais pretendem solucionar \textbf{problemas}.
As intenções pedagógicas devem promover a \textbf{mobilização} de processos cognitivos que permitam aos estudantes a plena interação com os objetos de conhecimento, transformando a sala de \textbf{aula} em um ambiente motivador e estimulante, que \textbf{desperte} as habilidades que irão ser requeridas em situações cotidianas ao longo da vida.
Tendo isso como base, é necessário criar possibilidades para \textbf{abordagens} \textbf{multidisciplinares}, incluindo temáticas orientadas para o pleno exercício da cidadania, ressaltando a relevância da \textbf{Ciência}, Tecnologia, Sociedade e Ambiente ( CTSA ) para a busca de soluções aos impactos \textbf{resultantes} do modo de vida \textbf{atual}.
A adoção de metodologias diversificadas, em conjunto com a utilização de dispositivos tecnológicos como alternativa ao \textbf{método} tradicional de ensino, abre novas possibilidades de cenários de aprendizagem por dinamizar o processo de ensino.
Quando inseridas como possibilidades de \textbf{abordagens}, as Metodologias Ativas, a exemplo da aprendizagem \textbf{baseada} em projetos e estudos de casos, podem \textbf{auxiliar} o professor na mediação do processo formativo dos estudantes.
Como referencial, as metodologias ativas adotam a participação ativa do estudante na condução do seu processo de aprendizagem, possibilitando o \textbf{despertar} do protagonismo, tornando, assim, a experiência de aprender mais estimulante e significativa.
As situações em que se desenvolvem as \textbf{aulas} sob o aspecto dinâmico podem facilitar o engajamento dos estudantes para compreender algo \textbf{relevante} ao seu cotidiano, em articulação com os Temas \textbf{Contemporâneos} \textbf{Transversais}, além de favorecer a interação entre os pares, entre professores-estudantes e entre estudantes-objetos de conhecimento, \textbf{mobilizando} emoções que \textbf{auxiliam} a \textbf{compreensão} de novos conceitos ou ressignificando os saberes já existentes.
Além das metodologias ativas no ensino de Química, Física e Biologia, há a necessidade de uma reflexão quanto à importância da implantação de atividades experimentais como metodologia constante na prática docente no Ensino Médio.
Para \textbf{tanto}, torna -se \textbf{relevante} a valorização de práticas experimentais alternativas ( a citar, uma \textbf{aula} de campo, nas imediações da escola, ou a utilização das áreas comuns da escola, como a cozinha escolar, ou mesmo a criação de hortas, composteiras e biodigestores que possam aproveitar a matéria orgânica \textbf{resultante} da produção da merenda escolar ), haja vista a diversidade social, \textbf{econômica} e ambiental na qual cada estabelecimento de ensino está inserido, considerando os diversos fatores que viabilizam ou não esse tipo de metodologia.
Dentre as possibilidades, importa também ressaltar que as práticas experimentais e investigativas como \textbf{abordagem} pedagógica são alternativas de intervenção que possibilitam ao estudante aguçar sua curiosidade natural por abordar fenômenos correlatos às situações vivenciadas no seu cotidiano.
A operacionalização de tais atividades, quando planejadas de acordo com o contexto no qual a escola se insere, favorece a construção do pensamento científico e uma melhor \textbf{compreensão} dos conceitos estruturantes das \textbf{ciências} e deve desenvolver a \textbf{capacidade} de argumentação \textbf{crítica} para o seu \textbf{posicionamento} como cidadão, respaldado em informações confiáveis com senso ético e de responsabilidade.
Desse modo, envolver os jovens em situações práticas, por meio do “ aprender fazendo ”, favorece a \textbf{sistematização} dos procedimentos necessários à experimentação, bem como a aproximação aos objetos de conhecimento, diminuindo assim o estranhamento a alguns temas ou conceitos dos componentes das Ciências da Natureza e suas Tecnologias.
Nesse cenário, o professor ( enquanto mediador ) que ensina Ciências pode propor a criação de jogos, dramatizações, composições musicais, com ou sem a utilização de ferramentas \textbf{digitais}, de forma que \textbf{instigue} a curiosidade do estudante e o incentive não apenas a entender os fenômenos, mas que os motive a ir em busca de novas descobertas, dotando -o de ferramentas que respaldam a segurança e a autoestima para que se enxergue também como \textbf{cientista} em \textbf{potencial}.
É importante ressaltar, ainda, a necessidade de compreender e acompanhar os anseios, as expectativas e as necessidades desse jovem em formação e buscar conectar esses anseios às novas necessidades da sociedade \textbf{contemporânea}.
Ao orientar uma atividade, o professor media a construção coletiva do conhecimento e o estudante reconhece que, enquanto realiza a ação, ele se responsabiliza também pela sua aprendizagem.
Integrar essa proposta ao currículo de \textbf{Ciências} da Natureza e suas Tecnologias obstina -se à superação do modelo \textbf{teórico} e conteudista tradicionalmente utilizado nas \textbf{aulas}.
A variedade de \textbf{abordagens} modifica a lógica naturalizada, na qual se observa o professor como o centro do processo educativo, e realoca o estudante à posição de \textbf{autor} protagonista da sua aprendizagem.
Uma \textbf{abordagem} dinâmica favorece a elasticidade cognitiva, que oportuniza uma melhor adaptação aos \textbf{eventos} sociais, cotidianos e estimula uma melhor resposta às situações que surgirão nas esferas pessoais ou profissionais do estudante, uma vez que utiliza uma diversidade de estímulos capazes de \textbf{aprimorar} habilidades já existentes ou de favorecer o \textbf{surgimento} de novas habilidades.
Uma aprendizagem significativa requer acima de tudo a intencionalidade do docente em inovar e amplificar as experiências de aprendizagens a fim de conduzir o estudante a um modelo de formação que \textbf{demande} um maior engajamento com os objetos de conhecimento. \textbf{Abordagens} dinâmicas e inovadoras devem ser o arcabouço para uma nova postura no ensino das Ciências, visando garantir uma preparação para que os jovens façam bom uso ao manipular os conhecimentos e que, assim, possam intervir de \textbf{maneira} positiva, justa e responsável na realidade em que estão inseridos.
Além de novas \textbf{abordagens}, a escola deve ser a responsável por organizar e interpretar os diversos conhecimentos que os estudantes trazem, principalmente quando \textbf{lembramos} que, enquanto cidadãos, precisam atuar socialmente, buscando selecionar, interpretar e \textbf{divulgar} diversas informações.
No processo de ensino e aprendizagem das Ciências da Natureza e suas Tecnologias, é necessário \textbf{mobilizar} os estudantes para a importância do valor de \textbf{aproximar} -se do mundo em que vivem, de questionarem sobre sua estrutura e sua natureza, de descobrirem o interesse em fazer perguntas e procurar e elaborar possíveis respostas para as situações e \textbf{problemas} reais enfrentados nas suas comunidades, confrontando concepções alternativas, reconhecendo e valorizando os diferentes saberes existentes ( além do eurocêntrico ), como a sabedoria popular.
As \textbf{abordagens} e procedimentos devem ser capazes de conduzir os estudantes a uma aprendizagem em Ciências, que se obstine a uma formação que contemple as dimensões ética, política, social, ambiental e tecnológica, oportunizando escolhas responsáveis para a melhoria da qualidade de vida individual e coletiva.
Assim, é fundamental que o professor seja o articulador que media os processos de ensino e aprendizagem, provocando, nas \textbf{aulas}, situações que estimulem o protagonismo dos estudantes no \textbf{decorrer} do processo formativo.
Como alternativa para otimizar as \textbf{aulas} e a produção de materiais didáticos e atividades diversas, \textbf{sugere} -se, dentre outros, a utilização dos seguintes recursos pedagógicos disponíveis na internet : 1.
ChemSketch : pacote de desenho que permite desenhar estruturas químicas, disponível em :. 2.
ChemDraw : editor de moléculas, disponível em :. 3.
Jmol : software de computador para modelagem molecular de estruturas químicas em 3 dimensões, disponível em :. 4.
GenChemLab : \textbf{aplicativo} \textbf{baseado} em OpenGL, que pretende simular vários exercícios laboratoriais de química geral, disponível em :. 5.
BkChem : editor de moléculas 2D gratuito escrito em Python por Beda Kosata, disponível em :. 6.
Mol View : consiste em duas partes principais, um editor de fórmula estrutural e um visualizador de modelo 3D, disponível em :. 7.
Avogadro : editor e visualizador de moléculas projetado para uso em plataforma cruzada em química computacional, modelagem molecular, bioinformática, \textbf{ciência} de materiais e áreas relacionadas, disponível em :. 8.
Gaussian : pacote de software de química computacional de propósito geral lançado inicialmente em 1970 por John Pople e seu grupo de pesquisa na Carnegie Mellon University como Gaussian 70, disponível em :. 9.
GaussView : visualizador de moléculas, disponível em :. 10.
PhetColorado : \textbf{Simulações} gratuitas de \textbf{ciências} e matemática para o ensino de \textbf{tópicos} STEM, incluindo física, química, \textbf{biologia} e matemática, da University of Colorado Boulder, disponível em :. 11.
LabVirt : O Laboratório Didático \textbf{Virtual} é uma iniciativa da Universidade de São Paulo - USP, atualmente coordenado pela Faculdade de Educação.
Nele você vai encontrar \textbf{simulações} feitas pela equipe do LabVirt a partir de roteiros de estudantes de ensino \textbf{médio} das escolas da rede pública ; links para \textbf{simulações} e sites interessantes encontrados na Internet ; exemplos de projetos na \textbf{seção}"projetos educacionais"e respostas de especialistas para questões enviadas através do site :. 12.
Revista Qnesc : A Revista Química Nova na Escola ( QNEsc ), com uma periodicidade trimestral, propõe -se a \textbf{subsidiar} o trabalho, a formação e a atualização da comunidade do Ensino de Química brasileiro.
QNEsc integra -se à linha editorial da Sociedade Brasileira de Química, que \textbf{publica} também a revista Química Nova e o Journal of the Brazillian Chemical Society.
Química Nova na Escola é um espaço aberto ao educador, suscitando \textbf{debates} e reflexões sobre o ensino e a aprendizagem de química.
Assim, contribui para a \textbf{tarefa} fundamental de formar verdadeiros cidadãos.
Nesse sentido, a Divisão de Ensino disponibiliza em seu portal, na íntegra, e de forma totalmente gratuita, todos os artigos \textbf{publicados} no formato PDF, disponível em :. 13.
Kahoot ( https://kahoot.com/ ), Quizziz ( https://quizizz.com/ ), QuizLet ( https://quizlet.com/pt-br ) : plataformas de Gamificação ; 14.
Socrative ( https://www.socrative.com/ ) : plataforma para realização de quiz, avaliações e atividades online por agendamento ou em tempo real ; 15.
Padlet ( https://padlet.com ) : plataforma para criar murais interativos e linhas do tempo ; 16.
AnswerGarden ( https://answergarden.ch/ ) : ferramenta para criação de chuva de ideias e nuvem de palavras ; 17.
Mentimeter ( https://www.mentimeter.com/ ) : plataforma para criação de votações e enquetes ; 18.
Jamboard ( https://jamboard.google.com/ ) : extensão do Google Chrome para estimular o trabalho colaborativo dos estudantes e que ajuda a criar murais interativos e linhas do tempo, compatível para utilizar em grupos e \textbf{aulas} on-line via Google Meet ou Zoom. 19.
Canva ( https://www.canva.com/ ) : plataforma para criação de materiais ( cards para mídias sociais, folhetos explicativos e slides ) ; 20.
PowToon ( https://www.powtoon.com/ ) e AdobeSpark ( https://spark.adobe.com/pt-BR/ ) : plataforma para criação de animações, slides e \textbf{vídeos} animados ; 21.
Scratch ( https://scratch.mit.edu/ ) : plataforma que utiliza a linguagem de programação para a criação e compartilhamento de jogos, \textbf{vídeos} e animações ; 22.
MakeBeliefComix ( https://www.makebeliefscomix.com/ ) : plataforma para criação de histórias em quadrinhos ; 23.
Mozaweb ( https://www.mozaweb.com/ ) : plataforma para aprendizagem interativa, com uma biblioteca de livros, cenas, \textbf{vídeos} e jogos 3D ; 24.
OBS Studio ( https://obsproject.com/ ), LOOM ( https://www.loom.com/ ) e OCam ( https://ocam.softonic.com.br/ ) : softwares para gravação de videoaulas e captura de tela ; 25.
Paperpile ( https://paperpile.com/app ) : \textbf{aplicativo} para organizar citações e referências de acordo com as normas técnicas. 26.
Tabela periódica - versão da IUPAC ( União Internacional de Química Pura e Aplicada ), disponível em :.
Destaca -se que as várias ferramentas listadas como sugestões acima podem ser trabalhadas em um planejamento, integrando diversas áreas de conhecimento.
É importante que os professores busquem se ambientar e se familiarizar com plataformas que podem, a princípio, parecer de difícil manuseio, mas que \textbf{agregam} valor ao trabalho pedagógico por favorecer o engajamento dos estudantes e também o desenvolvimento de habilidades que preconizam o \textbf{letramento} \textbf{digital} e tecnológico dos estudantes.
Ainda sobre o protagonismo dos estudantes, é interessante que o professor seja capaz de mediar a utilização das ferramentas para que os estudantes produzam os materiais, \textbf{vídeos} e animações e possam utilizá -las na construção do seu conhecimento, ou na criação e divulgação de conteúdos para mídias \textbf{digitais}.
As possibilidades são diversas e perpassam sobre a necessidade de inverter a perspectiva do modelo de \textbf{aula} no qual o professor detém todo o conhecimento.
A utilização de ferramentas \textbf{digitais}, familiares ao universo dos estudantes, pode possibilitar um trabalho em regime de colaboração no qual professores e estudantes ganham juntos, expertise.
Cabe esclarecer que as ferramentas \textbf{digitais} não são a única forma de promover engajamento e motivação, mas aqui são listadas como possibilidades \textbf{metodológicas} a serem acrescidas aos \textbf{métodos} já exitosos, tradicionais ou inovadores, utilizados no alcance da aprendizagem significativa.
Todos os recursos são louváveis quando se pretende guiar a aprendizagem pautada na intencionalidade da formação integral dos estudantes enquanto sujeitos de direitos.

% matched lemmas: abordagem, agregar, aplicativo, aprimorar, aproximar, atual, aula, autor, auxiliar, basear, biologia, capacidade, cientista, ciência, compreensão, contemporâneo, crítico, debate, decorrer, demandar, despertar, digital, divulgar, econômico, evento, instigar, lembrar, letramento, maneira, metodológico, mobilizar, mobilização, multidisciplinar, médio, método, posicionamento, potencial, problema, publicar, relevante, restringir, resultante, seção, simulação, sistematização, subsidiar, sugerir, surgimento, tanto, tarefa, tecnologia, teórico, transversal, tópico, virtual, vivar, vídeo
\end{textsample}
